\chapter{검증 체계}

\section{검증 수준과 필요성}

전력변환 시스템의 실패는 서로 다른 시간 규모와 원인에서 나타나므로 검증을 단순한 모델에서 자세한 모델로 올라가도록 구성했다(표 \ref{t:levels}). 한 수준의 통과는 그 수준의 가정 안에서만 유효하며, 위 수준의 완료로 넓혀 해석하지 않는다.

\Needspace{8\baselineskip}
{\footnotesize
\begin{longtable}{L{\dimexpr 0.1746\linewidth-2\tabcolsep\relax}L{\dimexpr 0.2222\linewidth-2\tabcolsep\relax}L{\dimexpr 0.3651\linewidth-2\tabcolsep\relax}L{\dimexpr 0.1270\linewidth-2\tabcolsep\relax}L{\dimexpr 0.1111\linewidth-2\tabcolsep\relax}}
\caption{검증 수준, 필요한 이유와 현재 상태}\label{t:levels}\\
\toprule
\tabhead \textbf{수준} & \textbf{모델·도구} & \textbf{필요한 이유} & \textbf{근거} & \textbf{상태} \\
\midrule
\endfirsthead
\toprule
\tabhead \textbf{수준} & \textbf{모델·도구} & \textbf{필요한 이유} & \textbf{근거} & \textbf{상태} \\
\midrule
\endhead
\bottomrule
\endlastfoot
① 해석·설계 추정 & 평형점, 규격 대조, 데이터시트 손실 규칙 & 정격 분리·보장 하한의 성립, 같은 규칙의 손실 비교 & E3, E4, E15, E18, E25, E31, E36, E41 & 완료 \\
①′ 부품 데이터시트 대조 & 모의 스트레스 대 데이터시트 표·본문 값, 미리 정한 여유 규칙 & 모델 스트레스를 실제 부품이 견디는지 확인 & E34, E37–E40 & 완료(조건부 포함) \\
② 평균모델 & Stage B 6상태 모델, 기준 구조, 창 조정 모듈, 보호 상태기계, 조정형 Stage C, V\textsubscript{A} 노드 & 부하·공차·지연·고장 조합 시험으로 제어기·정격·순서 결정 & E2, E5–E9, E13, E14, E16, E17, E19–E23, E26, E29, E38–E40 & 완료 \\
③ 독립 구현 & MATLAB 재구현, ngspice & 코드 오류와 설계 특성 구별 & E2, E11, E22 & 완료 \\
④ 스위칭 수준 모델 & Stage A 8위상(위상 셰딩 포함), Stage B(5 ns), Stage C 직류변압기(레그 출력용량), DAB 정상상태(데드타임·출력용량) & 리플, 데드타임, DAB 위상, 센서 오차, 영전압 전이 등 평균모델에 없는 현상 확인 & E1, E10, E12, E16, E21, E22, E24, E26–E28, E30, E32, E33, E35 & 완료 \\
⑤ 소자 수준 회로 & Simulink·Simscape Electrical & 방정식과 근사를 독립 회로 해석기로 점검, 세 단 연결, 결선·게이트 사양 확보 & TB01–TB33 & 완료(대표 사례) \\
⑥ HIL·시제품 & 실시간 시뮬레이터, 축소전력 시제품 & 실제 지연, 잡음, 열, 전자기 간섭 & — & 후속 단계 \\
\end{longtable}
}

\section{판정 원칙}

\begin{itemize}
\item 판정 기준은 시험 전에 고정하고(3장), 결과를 본 뒤에 바꾸지 않는다. 기준을 좁히거나 넓힐 때는 새 시험 번호로 다시 판정한다(예: TB10 회로 기준 2배 → TB11 1.5배). 기준을 아주 조금 넘은 경우도 불통과로 보고한다(예: TB20 사례 D, 1.01 A 대 1.0 A). 결과를 본 뒤 더한 진단은 판정이 아니다. 기록 순서의 한계는 18.2절에 적었다.
\item 모델은 실제 상태를 강제로 자르지 않는다. 제어기는 지령만 제한하므로 과부하나 지연으로 생기는 정격 위반이 그대로 드러난다.
\item 실패한 사례는 지우지 않고 회귀 시험으로 남긴다. 모델 사이의 차이(에지 양자화, 추가 소자, 해석기 특성)는 결과와 함께 명시한다.
\item 소자 수준 회로 모델은 검증 단계마다 새 테스트벤치를 만들어 남기고, 같은 조건의 검증 모델 결과를 기준값으로 둔다. 회로에서 새 현상이 드러나면 검증 모델로 먼저 재현한 뒤 설계를 고친다(예: TB14, TB17).
\end{itemize}

\section{검토 과정}

검증 결과는 아홉 차례의 외부 검토·재검토를 거쳤다. 지적 사항은 저장 결과 대조, 모델·소스 대조, 재현 시험, 문헌 검토 등으로 확인했고, 계산이 필요한 지적은 판정 기준을 실행 전에 정해 다시 판정했다. 재현되지 않은 차이와 미검증 사항은 한계로 따로 기록했다(TB33 C33, 18.4절). 검토별 지적과 대응은 표 \ref{t:reviews}에 정리했다. 2026-10-29의 6인 검토에서는 검토자 여섯이 최종 수치·부품 정격·진행 기록·재현·참고문헌·공학적 근거를 나눠 보았고, 핵심 수치 재실행 10개 부분은 저장 결과와 같았다.
